% 334_Superposition_ohne_Zeit_De_ch.tex
% Dok. 334, FFGFT-Korpus, August 2026 — Kapitel-Datei
% Superposition ohne Zeit
% Was die Abwesenheit des Zeitoperators in der QM bedeutet
% und wie T~·m=1 das strukturell auflöst

% ============================================================
\section*{Vorbemerkung}

Dieses Dokument baut auf Dok.~307 (Zeit im Zustandsraum oder
daneben) und Dok.~333 (Superposition in FFGFT) auf, ohne deren
Inhalt zu wiederholen. Der spezifische Fokus hier: Was bedeutet
es konkret für die Interpretation von
\emph{Superposition, Verschränkung und Kollaps}, dass die
Standard-QM keinen Zeitoperator hat? Und was ändert sich
strukturell, wenn $\tilde{T}\cdot m=1$ jeder Mode eine eigene
Zeitskala gibt?

% ============================================================
\section{Die klassische Vorgeschichte — Instantanität als Erbe}

\subsection{Klassische Mechanik: Zeit läuft, Berechnungen sind instantan}

Die Instantanität ist kein QM-spezifisches Problem. Sie ist ein
Erbe der klassischen Physik, das die QM unverändert übernommen hat.

In der Newton'schen Mechanik wirkt Gravitation instantan über
beliebige Distanzen. Das Coulomb-Potential ist instantan. Die
Lagrange- und Hamilton-Mechanik beschreibt den Zustand
$(q(t), p(t))$ zu einem fixen $t$ — die Berechnung ist eine
zeitlose algebraische Auswertung über den Momentanzustand.
Niemand fragt, wie lange diese Berechnung „dauert", weil
Berechnungen in dieser Sprache keine Dauer haben. Die Zeit
$t$ läuft im Hintergrund als externer Parameter — die Physik
steckt in der zeitlosen Struktur der Gleichungen.

Das ist dieselbe Struktur wie in der QM. Pauli hat nur
explizit gemacht, was klassisch stillschweigend war.

\subsection{Die verborgene Inkonsistenz der QM}

Die QM erbt diese Instantanität — und fügt eine neue Spannung
hinzu:

\textbf{Der Zustand $\psi(t)$ ist instantan.} Er ist eine
zeitlose mathematische Größe zu einem fixen $t$ — keine Dauer,
kein Prozess.

\textbf{Die Auswertung $|\psi(t)|^2$ braucht Zeit.}
$|\psi(t)|^2$ ist eine Wahrscheinlichkeitsdichte — um sie zu
verifizieren braucht man viele Messungen, also echte
physikalische Zeit. Die Statistik, die den Zustand belegt,
verteilt sich über einen zeitlichen Prozess. Der Formalismus
aber beschreibt nur den Zustand zu einem Moment.

Das ist die stillschweigende Inkonsistenz: \emph{Der Formalismus
ist instantan, seine Verifikation ist zeitlich ausgedehnt.}
Die Zeit, die für die probabilistische Auswertung gebraucht
wird, steht außerhalb des Formalismus — sie ist weder
Operator noch Variable, sondern das, was der Experimentator
aufwendet, um den Zustand überhaupt zu kennen.

\subsection{Konsequenz für Superposition}

Superposition $c_1\ket{\psi_1}+c_2\ket{\psi_2}$ ist damit
doppelt zeitlos: erstens als Zustand (kein Zeitoperator,
kein intrinsisches „Wann"), zweitens als Aussage
(probabilistisch, braucht Zeit zur Verifikation, die aber
außerhalb des Formalismus liegt). Was als überlagert
„existiert", kann nur durch wiederholte Messungen über
echte Zeit belegt werden — und jede einzelne Messung kollabiert
den Zustand, also zerstört die Superposition.

In FFGFT: die Periode $\tilde{T}_i = 1/m_i$ gibt der Mode
eine intrinsische Zeitskala. Die Superposition ist nicht
zeitlos — sie hat eine geometrisch definierte Skala, auf der
sie stabil oder instabil ist.

% ============================================================
\section{Die Abwesenheit von Zeit in der QM — unterschätzte Konsequenz}

\subsection{Das Pauli-Theorem und seine Folgen}

Wolfgang Pauli zeigte 1933, dass es für einen Hamiltonoperator
mit nach unten beschränktem Spektrum keinen selbstadjungierten
Zeitoperator geben kann, der kanonisch konjugiert zu $\hat{H}$
wäre (Dok.~307, Dok.~306). Die Konsequenz war stillschweigend
aber weitreichend: Zeit wurde aus dem Hilbertraum
herausgehalten und als klassischer externer Parameter behandelt.

Das ist nicht nur eine technische Entscheidung. Es bedeutet:
\textbf{Zeit ist in der QM kein Observabler.} Man kann nicht
fragen „wann ist das System in Zustand $\psi$?" und eine
Antwort aus dem Formalismus erhalten — man muss $t$ von
außen vorgeben.

\subsection{Was in der Lehre selten gesagt wird}

In den meisten Lehrbüchern wird Zeit als Parameter eingeführt,
ohne zu betonen was das fehlen eines Zeitoperators bedeutet.
Es bedeutet konkret:

\begin{itemize}
  \item Superposition $c_1\ket{\psi_1}+c_2\ket{\psi_2}$ ist
    ein Zustand \emph{ohne Zeitstruktur} — es gibt keine
    intrinsische „Uhr" die sagt, wie lange das System in
    dieser Superposition ist.
  \item Verschränkung zwischen zwei Teilchen ist
    \emph{zeitlos} — der gemeinsame Zustand hat keinen
    eingebauten Begriff von Gleichzeitigkeit.
  \item Kollaps ist \emph{ohne Dauer} — er „passiert" im
    externen Parameter $t$, aber $t$ selbst ist nicht Teil
    des Kollapsprozesses.
\end{itemize}

Das ist nicht ein Versagen der QM — es ist eine strukturelle
Eigenschaft, die selten explizit gemacht wird.

% ============================================================
\section{Drei Konsequenzen im Detail}

\subsection{Superposition ist zeitlos}

Ein Elektron im Zustand
$\ket{\psi} = c_\uparrow\ket{\uparrow} + c_\downarrow\ket{\downarrow}$
hat keine intrinsische Zeitskala. Die Frage „wie lange ist
es in dieser Superposition?" ist im QM-Formalismus nicht
stellbar — man bräuchte einen Zeitoperator, der nicht
existiert. Was existiert, ist die Kohärenzzeit $T_2$ — aber
das ist eine phänomenologische Größe, die Wechselwirkung mit
der Umgebung beschreibt, nicht eine intrinsische Eigenschaft
des Zustands.

In FFGFT ist das anders: Jede Mode $\mathcal{T}_i=(n_i,l_i,j_i)$
trägt $\tilde{T}_i = 1/m_i$ als geometrische Größe. Ein
Zustand \emph{ist} seine Zeitskala — nicht als Dauer einer
Superposition, sondern als strukturelle Eigenschaft der Mode.

\subsection{Verschränkung ohne Gleichzeitigkeit}

Zwei verschränkte Teilchen teilen einen gemeinsamen Zustand
$\ket{\Psi} \in \mathcal{H}_1\otimes\mathcal{H}_2$. Dieser
gemeinsame Zustand hat keine eingebaute Zeitstruktur. Die
Frage nach der Gleichzeitigkeit einer Messung an beiden
Teilchen ist daher keine QM-Frage, sondern eine klassische
Frage über den externen Parameter $t$ — und damit eine
relativistische Frage (Relativität der Gleichzeitigkeit).

Das ist der tiefste Grund, warum die QM mit Lorentz-Invarianz
vereinbar ist: weil „Gleichzeitigkeit" außerhalb des
Formalismus steht und daher von der Lorentz-Gruppe
transformiert werden kann, ohne den QM-Formalismus zu berühren.

In FFGFT: Wenn zwei verschiedene Moden mit
$m_1 \neq m_2$ verschränkt wären, hätten sie verschiedene
$\tilde{T}_1 \neq \tilde{T}_2$. Eine modenübergreifende
Verschränkung erzeugt daher inkompatible Zeitskalen —
strukturell problematischer als in der QM, wo dieser Konflikt
nicht auftritt, weil $t$ außerhalb steht.

\subsection{Kollaps ohne Dauer}

Der Messprozess in der QM beschreibt einen Übergang
$\ket{\psi} \to \ket{\psi_k}$ (Projektion auf einen
Eigenzustand). Dieser Übergang hat \emph{keine Dauer} im
QM-Formalismus — er ist instantan im Parameter $t$. Wie lange
der Kollaps dauert, ist eine Frage die außerhalb der
Standard-QM liegt (sie wird in Theorien des Messprozesses,
der Dekohärenz, oder der Quanten-Gravitation diskutiert).

Das ist keine Inkonsistenz, sondern direkte Konsequenz der
Abwesenheit des Zeitoperators: ohne $\hat{t}$ gibt es keine
Observable „Kollapsdauer".

In FFGFT ist der Typ-I-Schritt (Messakt, Dok.~285/333)
nicht instantan im abstrakten Sinn, sondern ein geometrischer
Übergang — der Massenkreis $S^1_m$ wird zur Zeitachse
$\mathbb{R}_t$, die Windungszahl wird verworfen. Das ist
ein Informationsverlust, der eine natürliche „Richtung" hat.

% ============================================================
\section{Was $\tilde{T}\cdot m = 1$ strukturell ändert}

\subsection{Zeit als Eigenschaft der Mode, nicht als Parameter}

Die Grundrelation $\tilde{T}\cdot m = 1$ (Dok.~285, 306, 307)
setzt Zeit nicht als externen Parameter, sondern als
geometrische Größe jeder Mode: Die Periode $R_m = \hbar/(mc^2)$
ist intrinsisch — sie gehört zum Teilchen wie seine Masse.

Das hat drei direkte Konsequenzen für die Punkte oben:

\textbf{Superposition:} Eine Superposition innerhalb einer
Mode — intermediäre Torsion (Dok.~174) — hat eine natürliche
Zeitskala $\tilde{T}_i$. Superposition \emph{über} verschiedene
Moden würde inkompatible Zeitskalen $\tilde{T}_i \neq \tilde{T}_j$
überlagern — das ist der strukturelle Grund, warum
modenübergreifende Superpositionen in FFGFT problematisch sind
(Dok.~333).

\textbf{Verschränkung:} Innerhalb einer Mode ist
Gleichzeitigkeit durch $\tilde{T}$ definiert. Über Moden
hinweg fehlt ein gemeinsamer Zeitbegriff — ähnlich wie in
der QM, aber aus einem anderen Grund: nicht weil $t$ extern
steht, sondern weil verschiedene Moden verschiedene
geometrische Zeitskalen tragen.

\textbf{Kollaps (Typ I):} Der Messakt (Dok.~285) hat eine
natürliche Richtung — verlustbehaftet von $S^1_m$ nach
$\mathbb{R}_t$ — und eine natürliche Asymmetrie: die Rückrichtung
$\mathbb{R}_t \to S^1_m$ erfordert die verworfene Windungszahl.
Der Zeitpfeil ist damit keine externe Zusatzannahme, sondern
Konsequenz der Geometrie des Messens.

\subsection{Pauli-Theorem und kompakte Zeit}

Das Pauli-Theorem gilt für einen \emph{offenen} Zeitparameter
konjugiert zu einem nach unten beschränkten Hamiltonian. Für
eine \emph{kompakte} Zeitdimension greift das Argument nicht
(Dok.~307): Die kompakte Zeit verhält sich wie eine
Winkelvariable, ihr Konjugiertes hat ein diskretes Spektrum
(Windungszahlen), und die Pauli-Bedingung ist nicht erfüllt.

$\tilde{T}\cdot m = 1$ realisiert genau diese Situation:
Zeit ist kompakt, das konjugierte Spektrum ist die diskrete
Massenleiter. Das Theorem schließt diesen Fall nicht aus —
es erzwingt ihn geradezu als den einzigen konsistenten Weg,
Zeit in den Hilbertraum zu integrieren.

% ============================================================
\section{Zusammenfassung und Abgrenzung}

\begin{center}
\begin{tabular}{p{3.5cm}p{4.5cm}p{4.5cm}}
\toprule
 & \textbf{Standard-QM} & \textbf{FFGFT ($\tilde{T}\cdot m=1$)} \\
\midrule
Zeit & absent: kein Operator, kein Hilbertraum-Element; externer
  Parameter & geometrisch: $\tilde{T}_i=1/m_i$, kompakt, modengebunden \\
Superposition & zeitlos, keine intrinsische Uhr & intermediäre
  Torsion innerhalb einer Mode; $\tilde{T}_i$ als Skala \\
Modenübergreifend & algebraisch erlaubt, kein Zeitkonflikt
  ($t$ extern) & inkompatible $\tilde{T}_i \neq \tilde{T}_j$,
  strukturell problematisch \\
Verschränkung & zeitlos, Gleichzeitigkeit extern (Lorentz) &
  modenspezifische $\tilde{T}$; modenübergreifend kein
  gemeinsamer Zeitbegriff \\
Kollaps & instantan in $t$, keine Dauer im Formalismus & Typ-I:
  gerichteter Informationsverlust, Zeitpfeil geometrisch \\
Pauli-Theorem & blockiert Zeitoperator für offenes $t$ &
  greift nicht für kompaktes $\tilde{T}$ (Winkelvariable) \\
\bottomrule
\end{tabular}
\end{center}

\medskip\noindent
Dieses Dokument macht nicht den Anspruch, die Standard-QM zu
ersetzen. Es zeigt, dass ihre strukturelle Eigenschaft —
Zeit als absent — selten explizit gemacht wird, und dass
$\tilde{T}\cdot m=1$ eine geometrische Alternative bietet,
die dieselben Phänomene (Superposition, Verschränkung,
Kollaps) aus einer anderen Grundstruktur beschreibt.

\medskip\noindent
\textbf{Verweise:} Dok.~174 (Superposition als Torsion),
Dok.~285 (Typ-I-Messakt), Dok.~306 (native Zeit-Energie-Reziprozität),
Dok.~307 (Zeit im Zustandsraum oder daneben),
Dok.~333 (Modenspezifität und Superposition).
